% !TeX root = ../papers/nc-p3.tex
\section{七次までの主元次数の比較と整な傾き}
\label{sec:rank-fifteen-smith-slopes}
\begin{lemma}\label{lem:rank-fifteen-smith-transfer}
局所整域 $B$ の分数体を $K$、剰余体を $k$ とし、行列の係数を
$B[[H]]$ に取る。$K[[H]]$ および $k[[H]]$ 上で、必要な本数までの
Smith標準形の $H$ 指数の列が等しいとする。
剰余体での単元主元を持ち上げて消去すると、各段の $H$ 除算を $B[[H]]$ 上で
厳密に実行でき、同じ本数の列を消去できる。
\end{lemma}
\begin{proof}
指数零の剰余主元を一つずつ使う。各主元は $B[[H]]$ の単元なので、残りの行列の
一般体上の指数列からも先頭の零が一つずつ消える。
目標本数に達する前に指数零を使い切ったなら、残りの一般体上の全成分は $H$ の倍数である。
$B\hookrightarrow K$ により全定数項は $B$ でも零である。全体を $H$ で割り、
以後の指数を一つずつ下げて続ける。
途中までの列に対する主張にも、最後に使用する指数を越える桁は要らない。
\end{proof}
$B=\mathbb Z[z]_{(3,m)}$、$k=\mathbb F_3[z]/m$ としてこの補題を使う。
比較用の一般体は常に $\mathbb Q(z)$ である。
基準の $m=z^2+z+2$ では第\ref{sec:rank-fifteen-three-residue-classes}節の
有理関数上の消去が、単元主元が尽きるたびに全定数項の恒等的な零を確認している。
従って基準剰余体で再計算した各主元の除算回数は一般体の指数列に等しい。
単に各行列の最大損失を比較するだけでは足りず、各指数の重複度を含めて照合する。

他の $m$ で未切断経路から順に比較する。前段で残る語は一般体上では
同じ飽和した語の格子の基底を与える。語を替えても $K[[H]]$ 上の基底変更は可逆である。
関係の末尾あるいは先頭に掛ける語の変更も可逆な行変換なので、次段の一般の指数列は不変である。
四中心の選択だけは最後の始点の四次の語を基準表示に揃え、
同じ自由座標を指定する。これで中心商の七次までの比較にも同じ議論が通る。
各段の準同型を剰余に移し、長さ8以後の残余の次元上界が第\ref{sec:rank-fifteen-three-residue-classes}節に
一致すれば、同節の三剰余類の補題で全次数の結論が得られる。

\begin{theorem}\label{thm:rank-fifteen-integral-depth}
標数零、剰余標数3の通常な階数15窓で $\nu(y)>0$ かつ $D=x^3/y$ が整とする。
全ての剰余傾きと全ての深さで、各始点で四つの中心四次式上自由階数256となり、
標準Koszul左右AS正則である。
\end{theorem}
\begin{proof}
\path{computations/verify_rank_fifteen_integral_depth.py} は、まず $\mathbb F_3(z)$ で
十五桁を取り、第\ref{sec:rank-fifteen-three-residue-classes}節の七次までの消去を行う。
主元の定数項は基準剰余体で単元とし、除算前の零は $\mathbb F_3(z)$ で確かめる。
剰余体で選んだ指数の各重複度は、同じ先行格子に対する一般体のものと一致する。
この消去の分母・主元の先頭分子が消えない他の表示では、同じ行列操作が同じ指数を与える。
従って補題\ref{lem:rank-fifteen-smith-transfer}により七次までの消去を混標数の局所環で行える。
七次の右積表の定数項を $\mathbb F_3(z)$ に取り、左右の積の照合では
九次に各文字六列を残す。八次から十三次まで $H$ 除算をせず、
各段の剰余上界は第\ref{sec:rank-fifteen-three-residue-classes}節と同じとなる。

この全計算で使用した単元の定数係数の分子・分母の平方自由な積の次数は12017である。
第\ref{sec:rank-fifteen-integral-slope}節の変換による最大公約元の反復で
残った次数96の多項式は、モニック化すると
\[
               (P-Q)(P+Q)(P^2+2PQ+2Q^2)
\]
のモニック化に等しい。$P/Q$ は同節の傾きである。
このため追加すべき剰余傾きは $1,-1$ および $T^2+2T+2$ の共役な二根に尽きる。

\path{computations/verify_rank_fifteen_smith_slopes.py} は、この順に
\[
 z^6+2z^4+2z^3+z^2+2,\qquad z^4+z^3+z^2+z+1,\qquad z^4+2z^3+2
\]
の根を取る。各多項式の既約性と所定の傾きを確認し、
対応する有限体で基準表示と同じすべての指数列を得る。
比較するのは未切断経路80組、中心方程式4組、最後の始点の基底と自由座標の照合4組、
七次までの中心商112組の合計200組である。
いずれの表示でもその後の八次から十三次の上界も一致する。
剰余表示を指定した $D$ へ持ち上げて $\nu(H)=\nu(y)/6>0$ とし、
同じ三剰余類の自由性判定を適用すれば全点で結論を得る。
\end{proof}
\begin{corollary}\label{cor:rank-fifteen-rational-integral-slopes}
$x,y\in\mathbb Q$ の通常な階数15窓で $D=x^3/y$ が3整なら、
各始点で自由階数256の標準Koszul左右AS正則窓である。
\end{corollary}
\begin{proof}
$x=0$ は既知なので $D\ne0$ とする。$\nu_3(D)>0$ は
系\ref{cor:rank-fifteen-positive-slope}で全点を覆う。
$D$ が単元なら $3\nu_3(x)=\nu_3(y)$ である。
$\nu_3(y)>0$ の場合は $\bar D=1,-1$ の両方が上の計算に含まれる。
負なら $\nu_3(y)<2\nu_3(x)$ により極の定理を使い、
零なら両座標が素体単元なので定理\ref{thm:rank-fifteen-unit-residues}が適用できる。
\end{proof}

\begin{theorem}\label{thm:rank-fifteen-pole-slope-annulus}
剰余標数3の通常な標数零窓で $\bar x=-1$、$0<2\nu(y)\le\nu(3)$ とする。
この窓は標準Koszul左右AS正則で、各始点で四中心四次式上自由階数256である。
この領域では $x^3/y$ は極を持つ。例えば $y^2=3$ の水平切片の $\bar x=-1$ の全点が含まれる。
\end{theorem}
\begin{proof}
$d^2=y/x^3$ とし、$(4+de)^3=e^2$ の単純剰余根 $\pm8$ を持ち上げる。
$U=1+de$, $r^2=U-1$, $T^2=x/(U+3)$ とすれば
\[
    U=1+r^2,\quad V=r^4 T^2,\quad x=(4+r^2)T^2,\quad y=r^4 T^6.
\]
$T$ は単元で $\bar T^2=-1$、$\nu(r)=\nu(y)/4$ である。
$A=(1-s^2)/(1+s^2)$, $B=2s/(1+s^2)$、$T=AB$ と選ぶ。
$T^2+1$ の分子は標数3で
$(s^4+2s^3+s+1)(s^4+s^3+2s+1)$ となる。
第一の既約因子の剰余根を指定し、$T$ の符号を選ぶ。
$T(1+s^2)^2-2s(1-s^2)$ の先頭と定数は単元で、指定根は持ち上がる。
$A^2+B^2=1$ より $K=\sqrt{1+r^2}$ に定数項1を選ぶと、
$n=rA,H/n=rB,H=r^2T$ は必要な平方和と積を持つ。

$P_\pm=\sqrt{2(K+1)\pm2Tr^2}$ の定数項はともに2である。
残りの根を $Q_\pm=r\sqrt{2/(K+1)\pm2T}$ とし、
根号内の平方根の定数項は順に $A+B,B-A$ を使う。
これらの定数項は単元である。
\path{computations/verify_rank_fifteen_pole_slope_annulus.py} は
まず $\mathbb Q(s)$ 上で十桁を使い、未切断経路の長さ1から5まで、
中心方程式、五次までの中心商の左右の積を計算する。
全ての除算前の定数係数は恒等的に零であり、全ての主元は上の四次因子で単元である。
従って第\ref{sec:rank-fifteen-overlap-annulus}節と同じく、
語の格子はこの範囲で自由であり、その階数は一般の下界に一致する。
次に $\mathbb F_3(s)$ 上で十七桁を使う。中心方程式の自由座標を含む
始点ごとに、四次の語基底と指定した座標が標数零のものと一致することも照合する。
これで同じ中心式を課した格子の積表を還元できる。
未切断経路の各長さの最大損失は $(0,1,1,2,2)$、
中心方程式はさらに一桁、五次までの中心商はさらに二桁を使い、八桁が残る。
仮定は $3\in r^8\mathcal O$ と同値なので、還元された五次の積表から
八桁で再開できる。長さ6から13で $(2,1,1,1,1,0,1,0)$ 桁を使い、最後に一桁残る。
商の次元上界は全始点で $(1+t+t^2+t^3)^4$ の係数となり、十三次で消える。
標数零の経路と中心の下界を合わせる同じ自由性・Koszul性の議論で結論を得る。
\end{proof}
